\documentclass[11pt]{article}

\usepackage[utf8]{inputenc}
\usepackage[T1]{fontenc}
\usepackage{lmodern}
\usepackage{geometry}
\usepackage{amsmath,amssymb,amsfonts,amsthm,mathtools}
\usepackage{bm}
\usepackage{enumitem}
\usepackage{microtype}
\usepackage{hyperref}
\usepackage{titlesec}
\usepackage{booktabs}
\usepackage{array}
\usepackage{setspace}
\usepackage{mathrsfs}
\usepackage{physics}

\geometry{margin=1in}
\hypersetup{colorlinks=true, linkcolor=black, urlcolor=blue, citecolor=black}
\setlist[enumerate]{topsep=0.4em,itemsep=0.35em}
\setlist[itemize]{topsep=0.3em,itemsep=0.25em}
\titleformat{\section}{\large\bfseries}{\thesection.}{0.5em}{}
\titleformat{\subsection}{\normalsize\bfseries}{\thesubsection.}{0.5em}{}
\setstretch{1.06}

\newcommand{\Aether}{\AE{}ther}
\newcommand{\Aflow}{\AE{}ther-flow}
\newcommand{\TheoryName}{\Aflow{} Interpretation of Relativity}
\newcommand{\TheoryProgram}{\Aether{} / \Aflow{} framework}
\newcommand{\CompletedMetric}{g^{\sharp}}
\newcommand{\RootCore}{\mathfrak R^{\mathrm{root,core}}}

\newtheorem{definition}{Definition}
\newtheorem{proposition}{Proposition}
\newtheorem{remark}{Remark}
\newtheorem{corollary}{Corollary}
\newtheorem{theorem}{Theorem}

\title{\textbf{The \TheoryName{}}\\[0.4em]
\large Primitive-Reservoir Observer-Reduction\\
\large Root Frontier Next Benchmark Criterion}
\author{}
\date{}

\begin{document}

\maketitle

\begin{abstract}
The active primitive-reservoir observer-reduction line has now moved below the former foundation frontier. The recorded root-side continuation already showed that the fixed foundation package is recovered canonically from a deeper root class, and the foundation-generating substrate root dynamics necessity / uniqueness theorem then spent the honest benchmark-facing compression move available there: once the fixed foundation package is required, the benchmark-relevant root core is forced up to root-faithful equivalence. After that result, another same-output root, foundation, or ground restatement no longer counts as the right next benchmark-facing manuscript, and neither does a reopening of stationary-reduction, stationary-Hessian, spectral, or selector layers.

This note records the resulting criterion. The next honest benchmark-facing manuscript must therefore do one of two sharper things: either derive or justify the foundation-generating substrate root dynamics from still more primitive substrate structure in a way that changes benchmark-facing status rather than merely relocating the unexplained layer deeper, or prove a genuinely smaller theorem-level collapse strictly inside the current benchmark-relevant root core \(\RootCore\). The benchmark boundary remains fixed throughout. The one operative metric is still exactly \(\CompletedMetric\), the ordinary matter route is unchanged, there is no second operative metric, the low-energy matter law is not enlarged, and no stronger promotion language is licensed. The positive result of the note is therefore routing and scope discipline at root scope.
\end{abstract}

\section{Introduction}

The active derivational program of \TheoryName{} is no longer at a stage where deeper continuation alone counts as benchmark-facing progress. The benchmark package is already fixed, the same-package observer-side remainder gate has already been settled on the active completed branch, and the recursive-depth safeguard now blocks same-output relays from counting as new frontier gains unless they also remove benchmark-relevant freedom or otherwise sharpen the routing burden.

At current scope, the foundation-generating substrate root dynamics necessity / uniqueness theorem is the manuscript that supplies that sharper gain. Its theorem-level content is not merely that the current root layer exists as one deeper realization of the same downstream package. It is that, once the fixed foundation package is demanded, the benchmark-relevant root core is no longer free. That is the live benchmark-facing frontier.

The present note records the routing consequence of that theorem. It does not reopen the root mathematics, and it does not claim a completed first-principles substrate derivation of gravity. It records only which kind of next manuscript would now count as an honest benchmark-facing gain below the current root frontier.

\paragraph{Framework and claim boundary.}
\TheoryProgram{} denotes the broader \Aether{} / \Aflow{} framework built on the claims that the \Aether{} is the underlying four-dimensional substrate of reality, the \Aflow{} is its intrinsic ordered motion, observed three-dimensional space is the local experiential slice of that deeper substrate, S-time is the experienced order of change arising from matter, light, and the \Aflow{}, and observed expansion is the three-dimensional appearance of deeper four-dimensional ordered motion. Within that framework, \TheoryName{} denotes the exact-closure theory in which the effective gravitational dynamics are adopted to be exactly Einsteinian with universal matter coupling. In the active sequence, \emph{adoption} means use of that established relativistic dynamics without claiming substrate derivation, while \emph{derivation} is reserved for a first-principles recovery from explicit substrate variables.

The present note stays strictly inside that boundary. It records only which kind of next manuscript would now count as an honest benchmark-facing gain at the current root frontier.

\section{Fixed Inputs}

\begin{definition}[Root-frontier fixed inputs]
The criterion recorded here uses the following fixed inputs.
\begin{enumerate}
    \item The benchmark exact-closure package, which fixes the public one-metric GR target of the repository.
    \item The post-bridge routing record, which already moved the live burden upstream from same-package observer-side closure to substrate derivation of that benchmark package.
    \item The former foundation-frontier compression theorem and the associated routing handoff, which moved the open burden below the benchmark-relevant foundation core.
    \item The recorded root-side continuation, which recovers the fixed foundation package from a deeper root class while preserving the same benchmark package.
    \item The foundation-generating substrate root dynamics necessity / uniqueness theorem, which proves that the benchmark-relevant root core \(\RootCore\) is forced up to root-faithful equivalence once the fixed foundation package is required.
\end{enumerate}
\end{definition}

\begin{remark}
The present note does not replace those manuscripts. It records the routing consequence of reading them together under the active benchmark safeguard against recursive depth drift.
\end{remark}

\section{Why the Root Necessity / Uniqueness Theorem Remains the Frontier}

\begin{definition}[Benchmark-neutral deeper root restatement]
A \emph{benchmark-neutral deeper root restatement} is a manuscript that preserves the same benchmark one-metric output, the same ordinary matter route, and the same settled observer-side closure verdict while merely relocating the unexplained substrate layer to a deeper root, foundation, or ground presentation without providing a new compression, necessity, uniqueness, or comparably sharp routing gain.
\end{definition}

\begin{proposition}[The root necessity / uniqueness theorem is the live frontier]
At the current root stage of the primitive-reservoir observer-reduction line, the foundation-generating substrate root dynamics necessity / uniqueness theorem remains the live benchmark-facing frontier.
\end{proposition}

\begin{proof}
That theorem removes benchmark-relevant freedom at root scope itself: once the fixed foundation package is required, the benchmark-relevant root core \(\RootCore\) is forced up to root-faithful equivalence. This is exactly the kind of necessity / uniqueness gain that changes benchmark-facing status under the routing rule. By contrast, a manuscript that only preserves the same benchmark package while pushing the unexplained stopping point deeper does not alter benchmark-facing status unless it adds a comparably sharp result. Therefore the live frontier remains the root necessity / uniqueness theorem rather than every later or parallel same-output relay recorded around it.
\end{proof}

\begin{corollary}[Status of the recorded root continuation]
The recorded root-to-foundation continuation is a legitimate active supporting manuscript, but under the current routing safeguard it is not by itself the default current benchmark-facing frontier.
\end{corollary}

\begin{proof}
The theorem records a deeper root class whose projected exact-shell output recovers the already fixed foundation package. That matters for continuity of the deeper chain. But because the induced benchmark-facing output remains the same and the benchmark-relevant internal root freedom is compressed only by the later necessity / uniqueness theorem, the recorded continuation does not on its own supersede the routing status of that frontier theorem.
\end{proof}

\section{The Next Benchmark Criterion}

\begin{definition}[Admissible next benchmark-facing root manuscript]
At the current root frontier, a manuscript counts as an admissible next benchmark-facing move only if it does at least one of the following:
\begin{enumerate}
    \item derives or justifies the foundation-generating substrate root dynamics from still more primitive substrate structure in a way that does more than merely relocate the unexplained layer deeper;
    \item proves a genuinely smaller theorem-level collapse strictly inside the benchmark-relevant root core \(\RootCore\);
    \item does either of the previous two while preserving the same benchmark one-metric package, the same ordinary matter route, and the same claim boundary.
\end{enumerate}
\end{definition}

\begin{theorem}[Next honest benchmark-facing task at root scope]
The next honest benchmark-facing manuscript at the current root frontier must either derive or justify the foundation-generating substrate root dynamics from still more primitive substrate structure in a benchmark-relevant way, or else prove a genuinely smaller collapse inside the benchmark-relevant root core. A manuscript that only repeats the current root, foundation, or ground data at a deeper level without such a new gain does not satisfy the criterion.
\end{theorem}

\begin{proof}
By Proposition 1, the live frontier is already the theorem that compresses the benchmark-relevant root freedom. Therefore a second manuscript at root scope must improve on that status rather than merely accompany it. A deeper-origin manuscript can do so only if it changes benchmark-facing status by more than depth alone. If it simply reproduces the same root-to-foundation package through another relay, the routing rule classifies it as benchmark neutral. The remaining honest alternatives are therefore exactly those stated in the definition: either a more primitive derivation that changes benchmark-facing status, or a sharper internal collapse of the current root core.
\end{proof}

\section{What the Next Move Should Not Be}

\begin{proposition}[Screened next moves]
At the current root frontier, none of the following is the default next benchmark-facing move:
\begin{enumerate}
    \item another same-output root, foundation, or ground restatement;
    \item a reopening of stationary-reduction, stationary-Hessian, spectral, or selector layers;
    \item a manuscript that introduces a second operative metric, enlarges the ordinary matter law, or uses stronger promotion language.
\end{enumerate}
\end{proposition}

\begin{proof}
Item (1) fails the preceding criterion because it does not directly address the current root-level burden. It revisits a recorded same-output relay rather than removing the remaining root-level freedom or proving a smaller internal core. Item (2) likewise does not directly address the live burden at current scope. It reopens lower or screened layers as if they were still frontier-defining rather than settling the present root-level task. Item (3) violates the fixed claim boundary of the benchmark package and therefore cannot count as a disciplined next step on the active line. The benchmark package remains the exact one-metric GR package already fixed by adoption, so any admissible next manuscript must preserve that boundary.
\end{proof}

\begin{remark}
This screening statement is not a denial that those lower-layer manuscripts exist or matter historically. It is a routing statement: they are not the default way to spend the next benchmark-facing move from the current root frontier.
\end{remark}

\section{Conclusion}

The positive result of this note is routing discipline at the current root frontier. The foundation-generating substrate root dynamics necessity / uniqueness theorem remains the live benchmark-facing gain because it already spends the honest internal compression move available at current root scope: the benchmark-relevant root core is forced once the fixed foundation package is demanded. The recorded root-side continuation remains active support, but under the recursive-depth safeguard it is not itself a new frontier gain.

The scope remains narrow and honest. The benchmark package is unchanged: the one operative metric is still exactly \(\CompletedMetric\), the ordinary matter route is unchanged, there is no second operative metric, and the low-energy matter law is not enlarged. Nothing here upgrades adoption to derivation or licenses stronger promotion language.

The next open burden is therefore clear. The next benchmark-facing manuscript should derive or justify the foundation-generating substrate root dynamics from still more primitive substrate structure in a way that changes benchmark-facing status, unless the sharper honest gain available instead is a genuinely smaller theorem-level collapse inside the benchmark-relevant root core. That is the clean next manuscript target at current scope.

\end{document}
